% !TeX root = ./slide.tex
% !TeX program = pdflatex
\documentclass[aspectratio=169,xcolor=table]{beamer}
\usetheme{Madrid}

% Definindo a cor vermelha institucional e aplicando ao tema
\definecolor{meuvermelho}{RGB}{180, 20, 20}
\usecolortheme[named=meuvermelho]{structure}

\usepackage[utf8]{inputenc}
\usepackage[T1]{fontenc}
\usepackage{lmodern}
\usepackage[portuguese]{babel}
\usepackage{graphicx}
\usepackage{booktabs}
\usepackage{multirow}
\usepackage{tikz}
\usetikzlibrary{shapes.geometric, arrows.meta, positioning, calc}
\usepackage{tcolorbox}
\usepackage{amsmath,amssymb}

% Paleta oficial do experimento
\definecolor{corHFSS}{HTML}{1B9E77}
\definecolor{corTidy}{HTML}{D95F02}
\definecolor{corAnalitico}{HTML}{7570B3}
\definecolor{corBloco}{HTML}{E7298A}
\definecolor{azulclaro}{HTML}{E6F2FF}
\definecolor{verdeclaro}{HTML}{EAFAEA}

% Moldura padronizada para as imagens dos 5 circuitos de RF
\newcommand{\fotocircuito}[2]{%
    \begin{tcolorbox}[
        colback=white, colframe=meuvermelho!80!black, boxrule=0.8pt, arc=3pt,
        boxsep=1.5pt, left=2pt, right=2pt, top=2pt, bottom=2pt
    ]
        \centering
        \includegraphics[width=0.98\textwidth, height=3.15cm, keepaspectratio]{Imagens/#1}\\[2pt]
        \tiny \textit{#2}
    \end{tcolorbox}%
}

% Moldura para as imagens de aplicações no Slide 3
\newcommand{\fotoapp}[3]{%
    \begin{tcolorbox}[
        colback=white, colframe=meuvermelho!80!black, boxrule=0.75pt, arc=3pt,
        boxsep=1.5pt, left=2pt, right=2pt, top=2pt, bottom=2pt
    ]
        \centering
        \includegraphics[width=0.98\textwidth, height=#2, keepaspectratio]{Imagens/#1}\\[2pt]
        \tiny \textbf{#3}
    \end{tcolorbox}%
}

% ==========================================
% CONFIGURAÇÃO DE TÍTULO E AUTORES
% ==========================================
\title[Comparação HFSS, Tidy3D e Analítico]{Comparação entre HFSS, Tidy3D e Modelo Analítico em Circuitos de Radiofrequência }
\subtitle{Proposta de Trabalho Final --- Planejamento Experimental em Blocos}

\author[PPGEE / UFMG]{
    Beatriz S. Campagnaro, Filipe C. Andrade e Izabela C. S. Barbosa
}

\institute[]{
    Disciplina: Planejamento e Análise de Experimentos \\
    Programa de Pós-Graduação em Engenharia Elétrica (PPGEE) --- UFMG
}
\date[10 de Novembro 2026]{10 de Novembro de 2026}

% Configuração das logos na parte inferior da capa
\titlegraphic{
    \begin{tikzpicture}[remember picture,overlay]
        % Logo da UFMG no canto inferior esquerdo
        \node[anchor=south west, yshift=0.5cm, xshift=1cm] at (current page.south west) {
            \includegraphics[height=1.4cm]{Imagens/logo_ufmg.jpg}
        };
        % Logo da Engenharia no canto inferior direito
        \node[anchor=south east, yshift=0.5cm, xshift=-1cm] at (current page.south east) {
            \includegraphics[height=1.6cm]{Imagens/logo_eng.png}
        };
    \end{tikzpicture}
}

\begin{document}

% ==========================================
% CAPA E ROTEIRO
% ==========================================
\begin{frame}
    \titlepage
\end{frame}

\begin{frame}{Roteiro da Apresentação}
    \tableofcontents
\end{frame}

% ==========================================
% SEÇÃO 1: CONTEXTUALIZAÇÃO E PROBLEMA DE PESQUISA
% ==========================================
\section{Importância da Simulação Eletromagnética e Problema}

\begin{frame}{Importância e Aplicações da Simulação Eletromagnética}
    \begin{columns}[c]
        % Coluna esquerda: apenas tópicos diretos
        \begin{column}{0.43\textwidth}
            \begin{block}{\footnotesize Aplicações e Relevância}
                \footnotesize
                \begin{itemize}\setlength{\itemsep}{6pt}
                    \item \textbf{Telecomunicações e RF:}
                    \begin{itemize}\scriptsize
                        \item Torres de transmissão, radares e enlaces 5G/6G.
                        \item Radiação de antenas em dispositivos móveis.
                    \end{itemize}
                    \item \textbf{Análise de Campos e Modos:}
                    \begin{itemize}\scriptsize
                        \item Visualização de modos ressonantes e distribuição de campo elétrico/magnético.
                        \item Compatibilidade eletromagnética (EMC).
                    \end{itemize}
                    \item \textbf{Prototipagem Virtual:}
                    \begin{itemize}\scriptsize
                        \item Previsão de acoplamentos e perdas antes da fabricação física.
                        \item Redução de custo e tempo de bancada.
                    \end{itemize}
                \end{itemize}
            \end{block}
        \end{column}

        % Coluna direita: imagens reais organizadas
        \begin{column}{0.54\textwidth}
            \begin{columns}[c]
                \begin{column}{0.44\textwidth}
                    \fotoapp{pngtree-electronic-radar-signal-antenna-tower-png-image_6938657.png}{4.8cm}{Torres de Comunicação e Radar}
                \end{column}
                \begin{column}{0.53\textwidth}
                    \fotoapp{antena_celular.png}{2.2cm}{Irradiação em Dispositivos Móveis}
                    \vspace{0.08cm}
                    \fotoapp{cavidade.jpeg}{2.2cm}{Distribuição de Campos e Modos}
                \end{column}
            \end{columns}
        \end{column}
    \end{columns}
\end{frame}

\begin{frame}{O Problema de Pesquisa: Evolução das Ferramentas}
    \centering
    % Linha do tempo / contraste visual entre os paradigmas
    \begin{tikzpicture}[scale=0.92, transform shape,
        pbox/.style={draw, line width=1pt, rounded corners=4pt, text width=3.9cm, minimum height=2.25cm, align=center, inner sep=5pt},
        arr/.style={-{Stealth[length=2.5mm]}, line width=1.2pt, draw=meuvermelho}
    ]
        \node[pbox, draw=corAnalitico!90!black, fill=corAnalitico!10] (p1) at (-5.05, 0.9) {
            \textcolor{corAnalitico!90!black}{\textbf{\small Teoria Clássica}}\\[3pt]
            \footnotesize \textbf{Modelo Analítico}\\[2pt]
            \scriptsize Dimensionamento rápido,\\mas despreza efeitos 3D
        };

        \node[pbox, draw=corHFSS!85!black, fill=corHFSS!10] (p2) at (0, 0.9) {
            \textcolor{corHFSS!85!black}{\textbf{\small Consolidado (1990)}}\\[3pt]
            \footnotesize \textbf{Ansys HFSS (FEM)}\\[2pt]
            \scriptsize Padrão-ouro da indústria\\para simulação de RF
        };

        \node[pbox, draw=corTidy!90!black, fill=corTidy!10] (p3) at (5.05, 0.9) {
            \textcolor{corTidy!90!black}{\textbf{\small Nova Geração (2021)}}\\[3pt]
            \footnotesize \textbf{Tidy3D (FDTD em GPU)}\\[2pt]
            \scriptsize Criado para óptica/fotônica;\\\textbf{suporte a RF no último ano}
        };

        \draw[arr] (p1.east) -- (p2.west);
        \draw[arr] (p2.east) -- (p3.west);
    \end{tikzpicture}

    \vspace{0.35cm}
    \begin{block}{Motivação e Pergunta do Experimento}
        \footnotesize
        \begin{itemize}\setlength{\itemsep}{5pt}
            \item \textbf{Motivação (Lacuna):} Recente expansão do \textbf{Tidy3D} (FDTD) para problemas de RF, coexistindo com o consolidado \textbf{HFSS} (FEM) e o \textbf{Modelo Analítico}.
            \item \textbf{Pergunta do Experimento:} \textit{``Os três métodos (\textbf{HFSS}, \textbf{Tidy3D} e \textbf{Modelo Analítico}) produzem resultados estatisticamente equivalentes na resolução de circuitos de RF?''}
        \end{itemize}
    \end{block}
\end{frame}

% ==========================================
% SEÇÃO 2: FERRAMENTAS AVALIADAS E OBJETIVO
% ==========================================
\section{Métodos Avaliados e Objetivo do Experimento}

\begin{frame}{Comparação entre as 3 Abordagens e Objetivo}
    \centering
    % Diagrama visual com fonte maior nos blocos e maior espaçamento até o objetivo
    \begin{tikzpicture}[scale=0.92, transform shape,
        card/.style={draw, line width=1.1pt, rounded corners=4pt, minimum width=4.45cm, minimum height=3.1cm, align=center, inner sep=5pt},
        arrow/.style={-{Stealth[length=2.6mm]}, line width=1.2pt, draw=meuvermelho}
    ]
        % Card 1: HFSS
        \node[card, draw=corHFSS!85!black, fill=corHFSS!8] (hfss) at (-4.9, 0.5) {
            \textcolor{corHFSS!85!black}{\textbf{\normalsize 1. Ansys HFSS}}\\[3pt]
            \footnotesize \textbf{Elementos Finitos (FEM)}\\[5pt]
            \scriptsize $\bullet$ Usado desde \textbf{1990}\\
            \scriptsize $\bullet$ Domínio da frequência\\
            \scriptsize $\bullet$ Malha tetraédrica adaptativa\\
            \scriptsize $\bullet$ Padrão de mercado em RF
        };

        % Card 2: Tidy3D
        \node[card, draw=corTidy!90!black, fill=corTidy!8] (tidy) at (0, 0.5) {
            \textcolor{corTidy!90!black}{\textbf{\normalsize 2. Tidy3D}}\\[3pt]
            \footnotesize \textbf{Diferenças Finitas (FDTD)}\\[5pt]
            \scriptsize $\bullet$ Criado em \textbf{2021} (óptica/fotônica)\\
            \scriptsize $\bullet$ Domínio do tempo via GPU\\
            \scriptsize $\bullet$ Malha cartesiana estruturada\\
            \scriptsize $\bullet$ \textbf{Último ano:} compatível com RF
        };

        % Card 3: Analítico
        \node[card, draw=corAnalitico!90!black, fill=corAnalitico!8] (ana) at (4.9, 0.5) {
            \textcolor{corAnalitico!90!black}{\textbf{\normalsize 3. Modelo Analítico}}\\[3pt]
            \footnotesize \textbf{Equações Fechadas}\\[5pt]
            \scriptsize $\bullet$ Literatura clássica de RF\\
            \scriptsize $\bullet$ Linhas de transmissão e modos\\
            \scriptsize $\bullet$ Cálculo teórico imediato\\
            \scriptsize $\bullet$ Referência inicial de projeto
        };

        % Caixa central de objetivo mais abaixo (maior distância dos quadros)
        \node[draw=meuvermelho, line width=1.1pt, fill=gray!6, rounded corners=4pt, minimum width=13.9cm, minimum height=1.45cm, align=center, inner sep=4pt] (obj) at (0, -2.55) {
            \small \textbf{Objetivo do Experimento: Testar a Equivalência entre os 3 Métodos ($a = 3$)}\\[3pt]
            \footnotesize Avaliar simultaneamente se \textbf{HFSS}, \textbf{Tidy3D} e \textbf{Analítico} são equivalentes ({$H_0: \tau_1 = \tau_2 = \tau_3 = 0$}),\\
            \footnotesize isolando os \textbf{5 circuitos} em blocos ($b = 5$) com \textbf{2 réplicas dielétricas} por bloco ($N = 30$).
        };

        \draw[arrow] (hfss.south) -- (hfss.south |- obj.north);
        \draw[arrow] (tidy.south) -- (tidy.south |- obj.north);
        \draw[arrow] (ana.south)  -- (ana.south  |- obj.north);
    \end{tikzpicture}
\end{frame}

% ==========================================
% SEÇÃO 3: OS 5 CIRCUITOS DE RF (BLOCOS)
% ==========================================
\section{Dispositivos de Teste: Os 5 Circuitos de RF}

\begin{frame}{Dispositivos Avaliados (Parte I): Guias e Estruturas Ressonantes}
    \begin{columns}[T]
        % Circuito 1
        \begin{column}{0.32\textwidth}
            \begin{block}{\scriptsize 1. Cavidade Ressonante Retangular}
                \fotocircuito{cavidade.jpg}{Estrutura metálica fechada 3D}
                \scriptsize
                \begin{itemize}\setlength{\itemsep}{2pt}
                    \item \textbf{Réplicas ($n=2$):} Ar e Rutila.
                    \item \textbf{Métrica ($M_{i1k}$):} Freq.\ $f_{101}$ ($\text{GHz}$).
                \end{itemize}
            \end{block}
        \end{column}

        % Circuito 2
        \begin{column}{0.32\textwidth}
            \begin{block}{\scriptsize 2. Linha de Transmissão Microfita}
                \fotocircuito{microfita.jpg}{Guia planar quase-TEM}
                \scriptsize
                \begin{itemize}\setlength{\itemsep}{2pt}
                    \item \textbf{Réplicas ($n=2$):} FR-4 e RO3010.
                    \item \textbf{Métrica ($M_{i2k}$):} Impedância $Z_0$ ($\Omega$).
                \end{itemize}
            \end{block}
        \end{column}

        % Circuito 3
        \begin{column}{0.32\textwidth}
            \begin{block}{\scriptsize 3. Divisor de Potência de Wilkinson}
                \fotocircuito{wilkinson.jpg}{Divisor de 3 portas ($-3\text{ dB}$)}
                \scriptsize
                \begin{itemize}\setlength{\itemsep}{2pt}
                    \item \textbf{Réplicas ($n=2$):} FR-4 e RO3010.
                    \item \textbf{Métrica ($M_{i3k}$):} Transmissão $|S_{21}|$ ($\text{dB}$).
                \end{itemize}
            \end{block}
        \end{column}
    \end{columns}
\end{frame}

\begin{frame}{Dispositivos Avaliados (Parte II): Radiação, Acoplamento e Padronização}
    \begin{columns}[T]
        % Circuito 4
        \begin{column}{0.32\textwidth}
            \begin{block}{\scriptsize 4. Antena de Microfita (\textit{Patch})}
                \fotocircuito{antena.jpg}{Elemento radiante planar}
                \scriptsize
                \begin{itemize}\setlength{\itemsep}{2pt}
                    \item \textbf{Réplicas ($n=2$):} FR-4 e RO3010.
                    \item \textbf{Métrica ($M_{i4k}$):} Freq.\ central $f_r$ ($\text{GHz}$).
                \end{itemize}
            \end{block}
        \end{column}

        % Circuito 5
        \begin{column}{0.32\textwidth}
            \begin{block}{\scriptsize 5. Acoplador Híbrido \textit{Branch-Line}}
                \fotocircuito{branchline.jpg}{Acoplador em quadratura ($90^\circ$)}
                \scriptsize
                \begin{itemize}\setlength{\itemsep}{2pt}
                    \item \textbf{Réplicas ($n=2$):} FR-4 e RO3010.
                    \item \textbf{Métrica ($M_{i5k}$):} Acoplamento $|S_{31}|$ ($\text{dB}$).
                \end{itemize}
            \end{block}
        \end{column}

        % Diagrama visual do isolamento, réplicas e normalização
        \begin{column}{0.32\textwidth}
            \begin{tcolorbox}[
                colback=azulclaro!45, colframe=meuvermelho, boxrule=0.9pt, arc=3pt,
                boxsep=3pt, left=3pt, right=3pt, top=4pt, bottom=4pt,
                title={\scriptsize \textbf{Blocos, Réplicas e Normalização}}
            ]
                \centering
                \vspace{0.15cm}
                \begin{tikzpicture}[
                    bnode/.style={draw=meuvermelho!80!black, rounded corners=3pt, fill=white, minimum width=3.45cm, minimum height=0.82cm, inner sep=2pt, align=center, font=\scriptsize},
                    arr/.style={-{Stealth[length=2mm]}, thick, draw=meuvermelho}
                ]
                    \node[bnode] (c1) at (0, 1.4) {\textbf{5 Blocos $\times$ 2 Dielétricos ($n=2$)}\\Amplia os dados ($N=30$ ensaios)};
                    \node[bnode, fill=verdeclaro] (c2) at (0, 0) {\textbf{Normalização Adimensional ($y_{ijk}$)}\\Mesma escala + $15\text{ GL}$ no resíduo};
                    \node[bnode, fill=corTidy!15] (c3) at (0, -1.4) {\textbf{Comparação Global ($a=3$)}\\HFSS vs.\ Tidy3D vs.\ Analítico};

                    \draw[arr] (c1) -- (c2);
                    \draw[arr] (c2) -- (c3);
                \end{tikzpicture}\par
                \vspace{0.15cm}
            \end{tcolorbox}
        \end{column}
    \end{columns}
\end{frame}

% ==========================================
% SEÇÃO 4: METODOLOGIA E CONTINUIDADE DO TRABALHO
% ==========================================
\section{Metodologia Planejada e Continuidade do Trabalho}

\begin{frame}{Metodologia Planejada: RCBD com Réplicas Dentro do Bloco}
    \begin{columns}[c]
        % Coluna 1: Estrutura do Experimento e Padronização
        \begin{column}{0.44\textwidth}
            \centering
            \begin{tcolorbox}[
                colback=white, colframe=meuvermelho, boxrule=0.8pt, arc=3pt,
                boxsep=2pt, left=2pt, right=2pt, top=3pt, bottom=3pt,
                title={\scriptsize \textbf{Estrutura de Controle ($N = 30$ Simulações)}}
            ]
                \centering
                \begin{tikzpicture}[scale=0.78, transform shape,
                    box/.style={draw=meuvermelho, thick, rounded corners=3pt, fill=gray!8, minimum width=2.5cm, minimum height=1.0cm, align=center, font=\scriptsize},
                    arrow/.style={-{Stealth[length=2.0mm]}, thick, draw=black!80}
                ]
                    \node[box, fill=verdeclaro] (proc) at (0,0) {\textbf{Resolução do}\\\textbf{Circuito de RF}};

                    \node[align=center, font=\tiny, draw=corHFSS!85!black, fill=corHFSS!12, rounded corners=2pt, inner sep=3pt] (fat) at (0, 1.65) {
                        \textbf{Fator de Interesse ($a=3$ Métodos):}\\
                        HFSS (FEM) $|$ Tidy3D (FDTD) $|$ Analítico
                    };

                    \node[align=center, font=\tiny, draw=corBloco!85!black, fill=corBloco!12, rounded corners=2pt, inner sep=3pt] (bloc) at (0, -1.65) {
                        \textbf{Fator de Bloco ($b=5$ Circuitos):}\\
                        Isola a variabilidade entre geometrias
                    };

                    \node[align=center, font=\tiny, draw=black!50, fill=gray!10, rounded corners=2pt, inner sep=2.5pt] (in) at (-2.85, 0) {
                        \textbf{Réplicas ($n=2$):}\\[1pt]
                        FR-4 / RO3010\\
                        \textit{(Ar / Rutila)}
                    };
                    \node[align=center, font=\tiny, draw=meuvermelho, fill=azulclaro, rounded corners=2pt, inner sep=3pt] (out) at (2.75, 0) {
                        \textbf{Resposta Norm.}\\[1pt]
                        $y_{ijk}$ (Adim.)
                    };

                    \draw[arrow] (in) -- (proc);
                    \draw[arrow] (proc) -- (out);
                    \draw[arrow] (fat) -- (proc);
                    \draw[arrow] (bloc) -- (proc);
                \end{tikzpicture}
            \end{tcolorbox}

            \vspace{0.10cm}
            \begin{tcolorbox}[colback=gray!8, colframe=meuvermelho!75!black, boxrule=0.6pt, arc=2pt, boxsep=2pt, left=2pt, right=2pt, top=2pt, bottom=2pt]
                \centering \tiny
                \textbf{Padronização Adimensional e Resíduo ($15\text{ GL}$):}\\[2pt]
                $y_{ijk} = \dfrac{M_{ijk} - M_{\text{ref},jk}}{M_{\text{ref},jk}}$ \quad $\Rightarrow$ \quad $y_{ijk} = \mu + \tau_i + \beta_j + (\tau\beta)_{ij} + \epsilon_{ijk}$
            \end{tcolorbox}
        \end{column}

        % Coluna 2: Matriz Experimental com Réplicas
        \begin{column}{0.54\textwidth}
            \centering
            \scriptsize \textbf{Matriz Experimental ($b = 5$ Blocos $\times$ $a = 3$ Métodos $\times$ $n = 2$ Dielétricos)}
            \vspace{0.1cm}

            \renewcommand{\arraystretch}{1.18}
            \resizebox{\textwidth}{!}{%
            \begin{tabular}{l l | c c c}
                \toprule
                \textbf{Bloco ($M_j$)} & \textbf{Dielétrico ($k$)} & \textcolor{corHFSS}{\textbf{HFSS ($y_{1jk}$)}} & \textcolor{corTidy}{\textbf{Tidy3D ($y_{2jk}$)}} & \textcolor{corAnalitico}{\textbf{Analítico ($y_{3jk}$)}} \\
                \midrule
                \multirow{2}{*}{\textbf{1. Cavidade ($f_{101}$)}} 
                    & $k_1$: Ar ($\varepsilon_r \approx 1$)           & $y_{111}$ & $y_{211}$ & $y_{311}$ \\
                    & $k_2$: Rutila ($\varepsilon_r \approx 82$)    & $y_{112}$ & $y_{212}$ & $y_{312}$ \\
                \midrule
                \rowcolor{gray!10}
                & $k_1$: FR-4 ($\varepsilon_r \approx 4{,}4$)    & $y_{121}$ & $y_{221}$ & $y_{321}$ \\
                \rowcolor{gray!10}
                \multirow{-2}{*}{\textbf{2. Microfita ($Z_0$)}}
                    & $k_2$: RO3010 ($\varepsilon_r = 10{,}2$)       & $y_{122}$ & $y_{222}$ & $y_{322}$ \\
                \midrule
                \multirow{2}{*}{\textbf{3. Wilkinson ($|S_{21}|$)}} 
                    & $k_1$: FR-4 ($\varepsilon_r \approx 4{,}4$)    & $y_{131}$ & $y_{231}$ & $y_{331}$ \\
                    & $k_2$: RO3010 ($\varepsilon_r = 10{,}2$)       & $y_{132}$ & $y_{232}$ & $y_{332}$ \\
                \midrule
                \rowcolor{gray!10}
                & $k_1$: FR-4 ($\varepsilon_r \approx 4{,}4$)    & $y_{141}$ & $y_{241}$ & $y_{341}$ \\
                \rowcolor{gray!10}
                \multirow{-2}{*}{\textbf{4. Patch ($f_r$)}}
                    & $k_2$: RO3010 ($\varepsilon_r = 10{,}2$)       & $y_{142}$ & $y_{242}$ & $y_{342}$ \\
                \midrule
                \multirow{2}{*}{\textbf{5. Branch-Line ($|S_{31}|$)}} 
                    & $k_1$: FR-4 ($\varepsilon_r \approx 4{,}4$)    & $y_{151}$ & $y_{251}$ & $y_{351}$ \\
                    & $k_2$: RO3010 ($\varepsilon_r = 10{,}2$)       & $y_{152}$ & $y_{252}$ & $y_{352}$ \\
                \midrule
                \rowcolor{azulclaro!60}
                \multicolumn{2}{l|}{\textbf{Média Adimensional ($\bar{y}_{i\cdot\cdot}$)}} & $\bar{y}_{1\cdot\cdot}$ & $\bar{y}_{2\cdot\cdot}$ & $\bar{y}_{3\cdot\cdot}$ \\
                \bottomrule
            \end{tabular}%
            }
        \end{column}
    \end{columns}
\end{frame}

\begin{frame}{Modelagem Estatística em Blocos e Hipóteses}
    \footnotesize
    Com as respostas normalizadas ($y_{ijk}$) e $n=2$ réplicas dielétricas por bloco ($N = 30$), adota-se o modelo RCBD com interação método-bloco para testar simultaneamente os \textbf{3 métodos}:

    \vspace{0.18cm}
    \centering
    % Fluxograma da decomposição do modelo e soma de quadrados
    \begin{tikzpicture}[scale=0.84, transform shape,
        eqbox/.style={draw=meuvermelho, line width=1pt, rounded corners=4pt, fill=gray!6, inner sep=5pt, align=center},
        partbox/.style={draw, line width=0.9pt, rounded corners=3pt, minimum width=3.15cm, minimum height=1.5cm, align=center, inner sep=3.5pt},
        arr/.style={-{Stealth[length=2.2mm]}, line width=1pt, draw=meuvermelho!85!black}
    ]
        % Equação principal no topo
        \node[eqbox, minimum width=13.6cm] (mod) at (0, 1.5) {
            \small $y_{ijk} = \mu + \tau_i + \beta_j + (\tau\beta)_{ij} + \epsilon_{ijk}, \qquad \epsilon_{ijk} \overset{\text{i.i.d.}}{\sim} \mathcal{N}(0, \sigma^2)$\\[3pt]
            \scriptsize $\tau_i$: Método ($i=1,2,3$) \;$|$\; $\beta_j$: Bloco/Circuito ($j=1,\dots,5$) \;$|$\; $(\tau\beta)_{ij}$: Interação \;$|$\; $k=1,2$: Réplicas (Dielétricos)
        };

        % Quatro caixas da partição da variância (Total = 29 GL)
        \node[partbox, draw=corHFSS!85!black, fill=verdeclaro] (ssm) at (-5.1, -0.5) {
            \scriptsize \textbf{Métodos ($\tau_i$)}\\[2pt]
            \footnotesize $SS_{\text{Métodos}}$ \scriptsize ($2\text{ GL}$)\\[2pt]
            \tiny Efeito global entre\\[-1pt]
            \tiny HFSS, Tidy3D e Analítico
        };

        \node[partbox, draw=corBloco!85!black, fill=corBloco!10] (ssb) at (-1.7, -0.5) {
            \scriptsize \textbf{Blocos ($\beta_j$)}\\[2pt]
            \footnotesize $SS_{\text{Blocos}}$ \scriptsize ($4\text{ GL}$)\\[2pt]
            \tiny Variabilidade isolada\\[-1pt]
            \tiny dos 5 circuitos de RF
        };

        \node[partbox, draw=corTidy!90!black, fill=corTidy!10] (ssint) at (1.7, -0.5) {
            \scriptsize \textbf{Interação ($(\tau\beta)_{ij}$)}\\[2pt]
            \footnotesize $SS_{\text{Interação}}$ \scriptsize ($8\text{ GL}$)\\[2pt]
            \tiny Sensibilidade do método\\[-1pt]
            \tiny à topologia do circuito
        };

        \node[partbox, draw=black!60, fill=gray!12] (sse) at (5.1, -0.5) {
            \scriptsize \textbf{Resíduo ($\epsilon_{ijk}$)}\\[2pt]
            \footnotesize $SS_{\text{Erro}}$ \scriptsize ($15\text{ GL}$)\\[2pt]
            \tiny Erro puro nas $n=2$\\[-1pt]
            \tiny réplicas ($ab(n-1)$)
        };

        \draw[arr] (ssm.north |- mod.south) -- (ssm.north);
        \draw[arr] (ssb.north |- mod.south) -- (ssb.north);
        \draw[arr] (ssint.north |- mod.south) -- (ssint.north);
        \draw[arr] (sse.north |- mod.south) -- (sse.north);

        % Hipóteses na base
        \node[draw=meuvermelho, line width=0.9pt, fill=azulclaro!50, rounded corners=4pt, minimum width=13.6cm, inner sep=5pt, align=center] (hip) at (0, -2.25) {
            \footnotesize \textbf{Hipóteses do Teste ANOVA Global ($F_0 = MS_{\text{Métodos}} / MS_{\text{Erro}}$):}\\[2pt]
            \footnotesize $H_0: \tau_{\text{HFSS}} = \tau_{\text{Tidy3D}} = \tau_{\text{Analítico}} = 0$ \quad (Os 3 métodos são equivalentes) \qquad vs. \qquad $H_1: \exists\, \tau_i \neq 0$ \quad (Ao menos um difere)
        };

        \draw[arr] (ssm.south) -- (ssm.south |- hip.north);
        \draw[arr] (sse.south) -- (sse.south |- hip.north);
    \end{tikzpicture}
\end{frame}

\begin{frame}{Continuidade do Trabalho e Resultados Esperados}
    \begin{table}[]
    \centering
    \footnotesize
    \renewcommand{\arraystretch}{1.45}
    \begin{tabular}{|c|l|l|}
    \hline
    \rowcolor{gray!15}
    \textbf{\textcolor{meuvermelho}{Etapa}} & \textbf{\textcolor{meuvermelho}{Conteúdo}} & \textbf{\textcolor{meuvermelho}{Descrição}} \\ \hline
    \textbf{\#1} & Modelagem ($N=30$) & Simulação dos 5 circuitos $\times$ 2 dielétricos nos 3 métodos \\ \hline
    \textbf{\#2} & Normalização e Blocos & Padronização adimensional ($y_{ijk}$) das métricas nos $b=5$ blocos \\ \hline
    \textbf{\#3} & Análise e Decisão & ANOVA em blocos ($15\text{ GL}$ no resíduo), teste $F_0$ e pós-teste \\ \hline
    \end{tabular}
    \end{table}

    \vspace{0.25cm}
    \begin{block}{Metodologia de Avaliação e Resultados Esperados}
        \footnotesize
        \begin{itemize}\setlength{\itemsep}{4pt}
            \item \textbf{Padronização e Réplicas ($N=30$):} Uso de 1 métrica normalizada ($y_{ijk}$) por circuito e $n=2$ dielétricos por bloco, garantindo escala adimensional única ($\bar{y}_{i\cdot\cdot}$) e $15\text{ GL}$ para o erro residual.
            \item \textbf{Resultado Esperado:} Isolar o efeito dos circuitos ($SS_{\text{Blocos}}$) e testar simultaneamente se os \textbf{3 métodos} (\textbf{HFSS}, \textbf{Tidy3D} e \textbf{Modelo Analítico}) são estatisticamente equivalentes entre si.
        \end{itemize}
    \end{block}
\end{frame}

\section{Referências}

\begin{frame}{Referências Bibliográficas}
    \footnotesize
    \begin{thebibliography}{99}
        \bibitem{pozar}
        D.~M.~Pozar, \emph{Microwave Engineering}, 4th~ed.\hskip 1em plus 0.5em minus 0.4em\relax Hoboken, NJ: John Wiley \& Sons, 2012. ISBN: 978-0-470-63155-3.

        \bibitem{hfss}
        Ansys Inc., ``Ansys HFSS: 3D High Frequency Electromagnetic Simulation Software,'' 2024. [Online]. Disponível em: \url{https://ansys.synopsys.com/products/electronics/ansys-hfss}.

        \bibitem{tidy3d}
        Flexcompute Inc., ``Tidy3D: Fast GPU-accelerated FDTD Electromagnetic Solver,'' 2024. [Online]. Disponível em: \url{https://tidy3d.simulation.cloud}.

        \bibitem{wiley_article}
        \emph{International Journal of Antennas and Propagation}, ``Electromagnetic and RF Circuit Analysis,'' Wiley/Hindawi. [Online]. DOI: \url{https://doi.org/10.1155/ijap/2733569}.

        \bibitem{leaky_wave}
        F.~Ghasemifard et~al., ``Silicon-based optical leaky wave antenna with narrow beam radiation,'' 2011. [Online]. Disponível em: \url{https://www.researchgate.net/publication/51193354}.
    \end{thebibliography}
\end{frame}

\end{document}
